\section{Aportes de MinCiencias a Tecnologías Digitales, Inteligencia Artificial y Tecnologías Cuánticas}

La contribución de MinCiencias a la soberanía digital durante el periodo analizado se concentra en tres frentes: la construcción de política pública para inteligencia artificial, la financiación de investigación aplicada en tecnologías digitales y la formación de capacidades científicas en territorios con brechas históricas. La Hoja de Ruta para el desarrollo y aplicación de la inteligencia artificial, presentada en febrero de 2024, organizó cinco entornos de trabajo: ética y gobernanza; educación, investigación e innovación; industrias innovadoras y emergentes; datos y organizaciones; y privacidad, ciberseguridad y defensa \parencite{minciencias_hoja_ruta_ia_pdf_2024,minciencias_informe_gestion_2024_pdf}. Este instrumento sirvió como insumo para la formulación del CONPES 4144, que fija una política nacional de IA con 106 acciones hasta 2030 y una inversión aproximada de 479.273 millones de pesos \parencite{conpes_4144_ia_pdf_2025}.

La dimensión participativa de esta agenda se evidencia en los Diálogos Regionales en Inteligencia Artificial. Entre abril y mayo de 2024, MinCiencias realizó encuentros en Tumaco, Bucaramanga, Villavicencio, Medellín, Yopal, Barranquilla, San Andrés, Ibagué, Zipaquirá, Duitama, Montería, Pereira y Cali, con 3.173 personas registradas. El objetivo fue diagnosticar conocimiento, apropiación, barreras y posibles soluciones territoriales basadas en IA, involucrando academia, comunidad científica, sector productivo, entidades públicas y sociedad civil \parencite{minciencias_dialogos_ia_pdf_2024}.

En financiación de proyectos, el programa ColombIA Inteligente constituye el núcleo más directamente asociado con IA, tecnologías aeroespaciales y tecnologías cuánticas. La Convocatoria 950 de 2024 financió cinco proyectos de I+D+i para soluciones mediante inteligencia artificial y ciencias del espacio, con 6.627 millones de pesos de MinCiencias y 1.889 millones de contrapartidas. Además, vinculó capacidades de talento mediante 40 jóvenes investigadores, 7 estudiantes de maestría y 3 estancias posdoctorales \parencite{minciencias_rendicion_2024_2025_pdf}. La Convocatoria 966 de 2025 amplió el alcance hacia ciencias y tecnologías cuánticas e inteligencia artificial; abrió el 25 de abril de 2025, cerró el 18 de junio de 2025 y tenía previsto publicar resultados definitivos en septiembre de 2025 \parencite{minciencias_rendicion_2024_2025_pdf}.

La agenda también incluye apropiación tecnológica temprana. Colombia Robótica desarrolló campamentos STEAM en 2024 con 1.200 niños, niñas y adolescentes y 80 docentes beneficiados, en Antioquia, Cundinamarca y Tumaco, con actividades en programación, robótica, diseño, realidad virtual e inteligencia artificial. En 2025, el programa entregó y adecuó 15 laboratorios STEAM en instituciones educativas de Tumaco, con equipamiento tecnológico, mobiliario especializado, kits de robótica y material didáctico \parencite{minciencias_rendicion_2024_2025_pdf}. Adicionalmente, en la respuesta a la emergencia del Catatumbo entre el 26 de enero y el 7 de febrero de 2025, MinCiencias reportó talleres de robótica, programación, electrónica maker, astronomía y cartografía participativa, con 430 niños, niñas, adolescentes y jóvenes beneficiados, 500 kits entregados y 26 talleres realizados \parencite{minciencias_rendicion_2024_2025_pdf}.

Finalmente, Orquídeas y la convocatoria de Investigación Fundamental amplían la base de capacidades científicas que sostiene la transformación digital. Orquídeas financió 107 proyectos en su versión 2023 por 20.611 millones de pesos y proyectó 119 proyectos en 2024 por 28.702 millones, vinculando doctoras y jóvenes investigadoras a proyectos de I+D+i con enfoque territorial y de género \parencite{minciencias_informe_gestion_2024_pdf}. La Convocatoria 937 de Investigación Fundamental incluyó proyectos con aprendizaje de máquina e inteligencia artificial, como investigación sobre aleaciones con imágenes multiescala y aprendizaje de máquina, y avances terapéuticos en péptidos antimicrobianos con IA \parencite{minciencias_informe_gestion_2024_pdf}.

